% !PRINTED 473
\weberpage{473}

Dies führt uns also darauf, die Permutationen von $m$ Ziffern
\[
0, 1, 2 \ldots m-1
\]
und ihre Eigenschaften im Allgemeinen genauer zu studiren.

Geht die Ziffer 0 in $a_0$, 1 in $a_1$, 2 in $a_2$ u.\,s.\,f., $m-1$ in
$a_{m-1}$ über, so muss

\weq{4}{a_0, a_1, a_2 \ldots a_{m-1}}

irgend eine andere Anordnung der Ziffern $0, 1, 2 \ldots m-1$ sein. Den
Uebergang von der einen zu der anderen Anordnung bezeichnen wir durch
ein Symbol, wie

\weq{5}{\pi_a = \begin{pmatrix} 0, & 1, & 2 & \ldots & m-1 \\ a_0, & a_1, & a_2 & \ldots & a_{m-1} \end{pmatrix}}

und nennen diesen Uebergang eine \gesperrt{Permutation}.

Die Anzahl aller möglichen Permutationen von $m$ Elementen ist
\[
\Pi(m) = 1 \,.\, 2 \,.\, 3 \ldots m,
\]
worunter auch die \gesperrt{identische Permutation}
\[
\begin{pmatrix} 0, & 1, & 2 & \ldots & m-1 \\ 0, & 1, & 2 & \ldots & m-1 \end{pmatrix},
\]
die jedes Element an seiner Stelle lässt, mitgezählt ist.

Das Symbol (5) ändert seine Bedeutung nicht, wenn man die einzelnen
Paare
\[
\begin{array}{ccccc}
0, & 1, & 2, & \ldots & m-1 \\
a_0, & a_1, & a_2, & \ldots & a_{m-1}
\end{array}
\]
anders anordnet. Ist also
\[
b_0, b_1, b_2 \ldots b_{m-1}
\]
irgend eine Anordnung der Ziffern $0, 1, 2 \ldots m-1$, so können wir
für (5) auch setzen:

\weq{6}{\pi_a = \begin{pmatrix} b_0, & b_1, & b_2 & \ldots & b_{m-1} \\ a_{b_0}, & a_{b_1}, & a_{b_2} & \ldots & a_{b_{m-1}} \end{pmatrix},}

was mit (5) völlig gleichbedeutend ist.

Hieraus lässt sich der sehr wichtige Begriff der \gesperrt{Zusammensetzung von
Permutationen} ableiten. Ersetzen wir einen Index $a$ durch einen Index
$b$, dann $b$ durch $c$, so ist der Erfolg derselbe, als wenn sogleich
$a$ durch $c$ ersetzt wird. Wenn wir also zwei Permutationen nach
einander ausführen, so ist das Ergebniss eine bestimmte Permutation,
die wir in folgender Weise darstellen können. Sind
%CONT>>
